\documentclass{beamer}

\usepackage[T2A]{fontenc}
\usepackage[utf8]{inputenc}
\usepackage[english,russian]{babel}
\input{settings.tex}


\title{$H_\infty$ управление}
\subtitle{Теория Управления, лекция 13}
\author{Сергей Савин}
\centering
\date{\mydate}



\begin{document}
	\maketitle
	

	
	
	\begin{frame}{Сингулярные значения}
		\begin{flushleft}
			
			Рассмотрим систему:
			
			\begin{equation}
				\begin{cases}
					\dot{\bo{x}} = \bo{A}\bo{x}+\bo{B}\bo{w}
					\\
					\bo{y} = \bo{C}\bo{x}
				\end{cases}
			\end{equation}
			
			Её передаточная функция: $G(s) = \bo{C} (s\bo{I} - \bo{A})^{-1} \bo{B}$. Частотная характеристика на частоте $\omega = \omega_0$:
			
			\begin{equation}
				G(j\omega_0) = \bo{C} (j\omega_0\bo{I} - \bo{A})^{-1} \bo{B}
			\end{equation}
			
			Спектральная норма матрицы частотной характеристики (её наибольшее сингулярное число) представляет собой максимально возможное усиление входного сигнала на частоте $\omega_0$:
			
			\begin{equation}
				||G(j\omega_0)||_2 =\ \bar{\sigma} (G(j \omega_0))
			\end{equation}
			
			
		\end{flushleft}
	\end{frame}
	
	
	\begin{frame}{Что такое H-infinity}
		\begin{flushleft}
			
			
			Термин «H-infinity» отсылает к H-бесконечности норме системы — наибольшему усилению входного сигнала, которое система может обеспечить на любой частоте:
			
			\begin{equation}
				||G||_{\infty} = \underset{\omega}{\sup} \ \bar{\sigma} (G(j \omega))
			\end{equation}
			
			H-infinity — это стратегия синтеза управления, которая гарантирует, что даже в наихудшем случае входной сигнал (возмущение) не будет усилен более чем на заданную величину $\gamma$.
			
		\end{flushleft}
	\end{frame}
	
	
	\begin{frame}{Условия H-infinity}
		\begin{flushleft}
			
			Для системы $G(s)$:
			
			\begin{equation}
				\begin{cases}
					\dot{\bo{x}} = \bo{A}\bo{x}+\bo{B}\bo{w}
					\\
					\bo{y} = \bo{C}\bo{x}
				\end{cases}
			\end{equation}
			
			Следующие условия эквивалентны (по определению H-бесконечности и сингулярных чисел):
			
			\begin{equation}
				||G||_\infty < \gamma \ \Longleftrightarrow \ G(j\omega)^* G(j\omega) \prec \gamma^2 I, \ \forall \omega
			\end{equation}
			
			Лемма Калмана — Якубовича — Попова (КЯП):
			
			\begin{equation}
				||G||_\infty < \gamma \ \Longleftrightarrow \ \int_0^\infty \left(\bo{y}\T \bo{y} - \gamma^2 \bo{w}\T \bo{w} \right) dt \leq 0
			\end{equation}
			
			\textcolor{mygray}{(при условии $\bo{x}(0) = 0$).}
			
		\end{flushleft}
	\end{frame}
	
	
	\begin{frame}{Критерий функции запаса}
		\begin{flushleft}
			
			Мы можем гарантировать $||G||_\infty < \gamma$, если найдём \emph{функцию запаса} $V( \bo{x}) > 0$ для всех $\bo{x} \neq \bo{0}$ и $V(\bo{0}) = 0$, такую что:
			
			\begin{equation}
				\dot V(\bo{x}) + \bo{y}\T \bo{y} - \gamma^2 \bo{w}\T \bo{w} \leq 0
			\end{equation}
			
			Выражение $\bo{y}\T \bo{y} - \gamma^2 \bo{w}\T \bo{w} $ называется \emph{интенсивностью поступления}.
			
		\end{flushleft}
	\end{frame}
	
	\begin{frame}{Критерий функции запаса — доказательство}
		\begin{flushleft}
			
			Проинтегрируем критерий от $t=0$ до $t = T$:
			
			\begin{equation}
				V(\bo{x}(T)) - V(\bo{x}(0)) + \int_{0}^{T}( \bo{y}\T \bo{y} - \gamma^2 \bo{w}\T \bo{w}) dt \leq 0
			\end{equation}
			
			Предположим нулевые начальные условия (рассматриваем чисто входно-выходное поведение), тогда $V(\bo{x}(0)) = 0$ и, учитывая, что $V$ положительно определена:
			
			\begin{equation}
				\int_{0}^{T} (\bo{y}\T \bo{y} - \gamma^2 \bo{w}\T \bo{w})dt  \leq  -V(\bo{x}(T)) \leq 0
			\end{equation}
			
			Поскольку это верно для всех $T$, переход $T \rightarrow \infty$ восстанавливает ранее приведённый критерий КЯП во временной области.
			
		\end{flushleft}
	\end{frame}
	
	
	\begin{frame}{Вывод условий в виде ЛМН, 1}
		\begin{flushleft}
			
			Рассмотрим функцию запаса вида $V( \bo{x}) =  \bo{x}\T  \bo{P}  \bo{x}$, где $ \bo{P} \succ 0$. Тогда критерий принимает вид:
			
			\begin{align}
				\dot{\bo{x}}\T  \bo{P}  \bo{x} + \bo{x}\T  \bo{P}  \dot{\bo{x}} + \bo{y}\T \bo{y} - \gamma^2 \bo{w}\T \bo{w} \leq 0
			\end{align}
			
			Подставим уравнения динамики:
			%
			\begin{align*}
				(\bo{A}\bo{x}+\bo{B}\bo{w})\T  \bo{P}  \bo{x} + \bo{x}\T  \bo{P} (\bo{A}\bo{x}+\bo{B}\bo{w}) + \bo{x}\T  \bo{C}\T \bo{C}\bo{x} - \gamma^2 \bo{w}\T \bo{w} \leq 0
			\end{align*}
			
			В матричной форме:
			%
			\begin{align}
				\begin{bmatrix}
					\bo{x} \\ \bo{w}
				\end{bmatrix}\T
				\begin{bmatrix}
					\bo{A}\T\bo{P}+\bo{P}\bo{A}+\bo{C}\T\bo{C} & \bo{P}\bo{B} \\
					\bo{B}\T\bo{P}  & -\gamma^2 \bo{I}
				\end{bmatrix}
				\begin{bmatrix}
					\bo{x} \\ \bo{w}
				\end{bmatrix}
				\leq 0
			\end{align}
			
			
			
		\end{flushleft}
	\end{frame}
	
	
	
	\begin{frame}{Вывод условий в виде ЛМН, 2}
		\begin{flushleft}
			
			Последнее неравенство должно выполняться для всех $\bo{x}$, $\bo{w}$, поэтому его можно переписать в виде ЛМН:
			%
			\begin{align}
				\begin{bmatrix}
					\bo{A}\T\bo{P}+\bo{P}\bo{A}+\bo{C}\T\bo{C} & \bo{P}\bo{B} \\
					\bo{B}\T\bo{P}  & -\gamma^2 \bo{I}
				\end{bmatrix}
				\prec 0
			\end{align}
			
			Это стандартное ЛМН для проверки H-бесконечности нормы системы.
			
		\end{flushleft}
	\end{frame}
	
	
	\begin{frame}{Синтез H-infinity управления, 1}
		\begin{flushleft}
			
			Рассмотрим систему:
			
			\begin{equation}
				\begin{cases}
					\dot{\bo{x}} = \bo{A}\bo{x}+\bo{B}_1\bo{u}+\bo{B}_2\bo{w}
					\\
					\bo{y} = \bo{C}\bo{x}
				\end{cases}
			\end{equation}
			
			Предположим линейное управление $\bo{u} = \bo{K}\bo{x}$, тогда первое уравнение принимает вид $\dot{\bo{x}} = (\bo{A}+\bo{B}_1\bo{K})\bo{x}+\bo{B}_2\bo{w}$. Подставим это в ранее выведенные условия:
			%
			\begin{align}
				\begin{bmatrix}
					(\bo{A}+\bo{B}_1\bo{K})\T\bo{P}+\bo{P}(\bo{A}+\bo{B}_1\bo{K})+\bo{C}\T\bo{C} & \bo{P}\bo{B}_2 \\
					\bo{B}_2\T\bo{P}  & -\gamma^2 \bo{I}
				\end{bmatrix}
				\prec 0
			\end{align}
			
			Теперь нужно преобразовать это в ЛМН (линейное относительно переменных решения).
			
		\end{flushleft}
	\end{frame}
	
	
	\begin{frame}{Синтез H-infinity управления, 2}
		\begin{flushleft}
			
			Разделим неравенство на $\gamma$ и определим $\bar{\bo{P}} = \frac{1}{\gamma} \bo{P}$:
			
			\begin{align}
				\begin{bmatrix}
					(\bo{A}+\bo{B}_1\bo{K})\T\bar{\bo{P}}+\bar{\bo{P}}(\bo{A}+\bo{B}_1\bo{K})+\frac{1}{\gamma}\bo{C}\T\bo{C} & \bar{\bo{P}}\bo{B}_2 \\
					\bo{B}_2\T\bar{\bo{P}}  & -\gamma \bo{I}
				\end{bmatrix}
				\prec 0
			\end{align}
			
			Выполним конгруэнтное преобразование с матрицей $\begin{bmatrix}
				\bo{Q} &\bo{0} \\
				\bo{0}  &  \bo{I}
			\end{bmatrix}$, где $\bo{Q} = \bar{\bo{P}}^{-1}$:
			
			\begin{align}
				\begin{bmatrix}
					\bo{Q}(\bo{A}+\bo{B}_1\bo{K})\T+(\bo{A}+\bo{B}_1\bo{K})\bo{Q}+\frac{1}{\gamma}\bo{Q}\bo{C}\T\bo{C}\bo{Q} & \bo{B}_2 \\
					\bo{B}_2\T  & -\gamma \bo{I}
				\end{bmatrix}
				\prec 0
			\end{align}
			
			Применим дополнение Шура, чтобы линеаризовать квадратичный член, и определим $\bo{L} = \bo{K}\bo{Q}$, получим ЛМН:
			%
			\begin{align}
				\begin{bmatrix}
					\bo{Q}\bo{A}\T + \bo{A}\bo{Q} + \bo{L}\T \bo{B}_1\T +\bo{B}_1 \bo{L} & \bo{B}_2 & \bo{Q}\bo{C}\T \\
					\bo{B}_2\T  & -\gamma \bo{I} & \bo{0} \\
					\bo{C}\bo{Q} & \bo{0} & -\gamma \bo{I}
				\end{bmatrix}
				\prec 0
			\end{align}
			
			
		\end{flushleft}
	\end{frame}
	
	
\begin{frame}[fragile]{Синтез H-infinity управления: пример на CVXPY}
		\begin{flushleft}
			{\small
			\begin{verbatim}
import cvxpy as cp

n,m,p = A.shape[0],B1.shape[1],B2.shape[1]
eps = 1e-6*np.eye(n)
Q = cp.Variable((n,n), PSD=True)
L = cp.Variable((m,n))
g = cp.Variable(nonneg=True)
M = cp.bmat([
[Q@A.T + A@Q + L.T@B1.T + B1@L,  B2,      Q@C.T],
[B2.T,                          -g*np.eye(p), 0],
[C@Q,              0,      -g*np.eye(C.shape[0])]
])
prob = cp.Problem(cp.Minimize(g), [Q >> eps, M << 0])
prob.solve()

K = L.value @ np.linalg.pinv(Q.value)
			\end{verbatim}}
		\end{flushleft}
	\end{frame}
	
	
	
\myqrframe
	
	
	
	
	\begin{frame}{Appendix A}
		\framesubtitle{Congruence transformation and definiteness}
		\begin{flushleft}
			
			Consider matrices $\bo{P} \succ 0$, and $\bo{V} \in \R^{n, n}$, the latter is full rank. We can prove that:
			
			\begin{equation}
				\bo{P} \succ 0 \implies \bo{V}^\top\bo{P}\bo{V} \succ 0
			\end{equation}
			
			Proof: $\bo{x}^\top\bo{V}^\top\bo{P}\bo{V}\bo{x} = \bo{z}^\top\bo{P}\bo{z}$, where $\bo{z} = \bo{V}\bo{x}$. Since $\bo{P} \succ 0$, $\bo{z}^\top\bo{P}\bo{z} \geq 0$, hence $\bo{x}^\top\bo{V}^\top\bo{P}\bo{V}\bo{x} \geq 0$. 
			
			\begin{definition}
				Congruence transformation preserves semi-definiteness: $\text{det}(\bo{V}) \neq 0, \ \bo{P} \succ 0 \implies \bo{V}^\top\bo{P}\bo{V} \succ 0$
			\end{definition}
			
			
		\end{flushleft}
	\end{frame}
	
	
	
	
\end{document}
